\begin{trapbox}
	\begin{description}[style=nextline, leftmargin=2.8em, labelsep=0.5em, itemsep=0.35em]
		\item[\textbf{\textcolor{CTrap}{易错点一（射影判定错误）}}] 错将斜线上任意点在平面内的投影当作射影，忽略射影必须与斜足相连。
		\item[\textbf{\textcolor{CTrap}{正确理解（正确做法）}}] 射影是斜线上一点（常取垂足）向平面作垂线，垂足与斜足连线段。如图，$PO\perp\alpha$，则$OA$即为斜线$PA$的射影。
		\item[\textbf{\textcolor{CTrap}{错因分析（概念模糊）}}] 对射影的定义掌握不准确，作图时忽略垂足位置。
	\end{description}

	\medskip
	\begin{description}[style=nextline, leftmargin=2.8em, labelsep=0.5em, itemsep=0.35em]
		\item[\textbf{\textcolor{CTrap}{易错点二（忽略平面条件）}}] 使用定理时未验证直线$l$是否在平面$\alpha$内，直接对空间任意直线使用。
		\item[\textbf{\textcolor{CTrap}{正确理解（确认平面条件）}}] 定理中$l$必须是平面$\alpha$内的直线。应用前需明确$l\subset\alpha$。
		\item[\textbf{\textcolor{CTrap}{错因分析（条件遗漏）}}] 记忆定理时只关注垂直关系，忽视直线位置限制。
	\end{description}

	\medskip
	\begin{description}[style=nextline, leftmargin=2.8em, labelsep=0.5em, itemsep=0.35em]
		\item[\textbf{\textcolor{CTrap}{易错点三（方向混淆）}}] 已知$l\perp PA$误以为可直接得$l\perp OA$，忘记需要$PO\perp l$或混淆正逆条件。
		\item[\textbf{\textcolor{CTrap}{正确理解（分清正逆）}}] 若$l\perp OA$则$l\perp PA$（正向）；若$l\perp PA$则$l\perp OA$（逆向）。两个方向均需利用$PO\perp l$（由$PO\perp\alpha$自动保证）。
		\item[\textbf{\textcolor{CTrap}{错因分析（死记硬背）}}] 未理解定理推导过程，只会机械套用，导致条件误用。
	\end{description}
\end{trapbox}
